\documentclass[11pt]{article}

\usepackage[a4paper,margin=1in]{geometry}
\usepackage{booktabs}
\usepackage{array}
\usepackage{longtable}
\usepackage{hyperref}
\usepackage{amsmath}
\usepackage{enumitem}
\usepackage{xcolor}

\hypersetup{
  colorlinks=true,
  linkcolor=blue,
  urlcolor=blue
}

\title{Notes on the Current Reeb-Space Sheet Tracking Analysis}
\author{}
\date{}

\begin{document}
\maketitle

\section{Purpose}

This document explains how the current Reeb-space tracking results were
inspected, what the reported metric summaries mean, what the exploratory
``continuity threshold'' means, and what additional implementation is needed
before these results should be used in a paper.

The goal is not yet to write the paper. The goal is to understand the current
outputs well enough to decide which results are concrete, which are only
heuristic, and which additional diagnostics are needed.

\section{Data and Files Inspected}

The implementation is located at:

\begin{verbatim}
/home/mohit/Desktop/postdoc/timeVaryingReebSpace/scripts/reeb-flow-visualizer
\end{verbatim}

The generated datasets inspected were:

\begin{verbatim}
/media/mohit/8tbh/postdoc/timeVaryingReebFeatures/MVK_s1
/media/mohit/8tbh/postdoc/timeVaryingReebFeatures/MVK_s2
/media/mohit/8tbh/postdoc/timeVaryingReebFeatures/stilbene
/media/mohit/8tbh/postdoc/timeVaryingReebFeatures/torusGrowing
/media/mohit/8tbh/postdoc/timeVaryingReebFeatures/torusMoving
\end{verbatim}

For each dataset, the most important files are:

\begin{verbatim}
sankey/unified_sankey_viewer/data.json
sankey/sheet_overlaps.json
compareSheetShapesCache/results/sheet_shape_matches.json
compareSheetShapesCache/results/sheet_shape_summary.json
compareSheetShapesCache/cache/global_bounds.json
sheetRendering/<step>/step_<id>.png
sheetFiberSurfaceImages/<step>/fiber_surface_images_manifest.json
\end{verbatim}

The relevant code files are:

\begin{verbatim}
stage_02_build_sankey_data.py
stage_03_compute_sheet_overlaps.py
compareSheetShapes/compare_sheet_shapes.py
unified_sankey_viewer/stage_07_unified_sankey_viewer.py
\end{verbatim}

\section{What Is Being Matched?}

For each timestep, the pipeline reads the RSI file and extracts the top
\(N\) sheets by sheet area. In the current settings, \(N=20\). For each
adjacent timestep pair \(T_i,T_{i+1}\), the shape comparison stage compares
every top sheet in \(T_i\) against every top sheet in \(T_{i+1}\). This gives up
to \(20 \times 20 = 400\) candidate sheet-to-sheet matches per adjacent pair.

The matching is not a one-to-one assignment. The output stores candidate scores
for all source-target pairs. The viewer or any downstream analysis can then
choose how to display or filter them.

\section{Current Matching Metrics}

The current shape comparison code computes the following metrics.

\subsection{Shape IoU}

The sheet polygons are rasterized into masks in the bivariate range plane. The
shape intersection-over-union is

\[
\operatorname{shapeIoU}(A,B)
=
\frac{|M_A \cap M_B|}{|M_A \cup M_B|},
\]

where \(M_A\) and \(M_B\) are the raster masks of the two sheets. This measures
how similar the sheet footprints are in range space.

\subsection{Support Jaccard}

The RSI files store regular domain vertices associated with each sheet. The
support Jaccard score is

\[
\operatorname{supportJaccard}(A,B)
=
\frac{|V_A \cap V_B|}{|V_A \cup V_B|},
\]

where \(V_A\) and \(V_B\) are the regular vertex sets of the two sheets. This is
a domain-overlap measure, not a range-shape measure.

Important caveat: if sheets have zero regular vertices, the score is zero. This
does not necessarily mean that the sheet has no geometric extent in range
space. It means the stored regular-vertex support gives no overlap evidence.

\subsection{Area Ratio}

\[
\operatorname{areaRatio}(A,B)
=
\frac{\min(\operatorname{area}(A),\operatorname{area}(B))}
       {\max(\operatorname{area}(A),\operatorname{area}(B))}.
\]

This measures whether the sheet areas are similar. It does not measure whether
the sheets occupy the same location in range space.

\subsection{Bounding-Box IoU}

The range-space bounding boxes of the two sheets are compared using standard
box intersection-over-union. This is a cheaper, coarser version of shape IoU.

\subsection{Centroid Similarity}

Let \(c_A\) and \(c_B\) be the range-space centroids of the two sheets. Let
\(D\) be the diagonal length of the global range-space bounds. The current code
uses

\[
\operatorname{centroidSimilarity}(A,B)
=
\exp\left(-\frac{\|c_A-c_B\|}{D}\right).
\]

This score is often high for many candidates because the exponential decays
slowly over the global range diagonal. Therefore, it is useful as a weak
supporting cue, but not very discriminative by itself.

\subsection{Combined Score}

The combined shape score is a weighted sum:

\[
\begin{aligned}
C(A,B) ={}&
0.40\,\operatorname{shapeIoU}(A,B) \\
&+ 0.30\,\operatorname{supportJaccard}(A,B) \\
&+ 0.15\,\operatorname{areaRatio}(A,B) \\
&+ 0.10\,\operatorname{bboxIoU}(A,B) \\
&+ 0.05\,\operatorname{centroidSimilarity}(A,B).
\end{aligned}
\]

These weights come from \texttt{SHAPE\_SCORE\_DEFAULT\_WEIGHTS} in
\texttt{common.py}.

\section{How I Judged Which Metrics Look Useful}

The previous summary said that shape IoU appears to be the most useful metric.
That statement was based on exploratory diagnostics, not on a formal user
study or ground-truth validation.

Four checks were used.

\subsection{Distribution}

A useful metric should have enough spread to distinguish good and poor
candidate matches. If a metric is almost always near zero, it may be too sparse
for standalone matching. If it is almost always near one, it may not
discriminate among candidates.

In the real datasets, support Jaccard is usually very small, while centroid
similarity is usually high. Shape IoU and bounding-box IoU have more useful
spread.

\subsection{Correlation with the Current Combined Score}

For all candidate sheet pairs, I computed how each metric correlates with the
current combined score. This tells us which metrics are most consistent with
the score currently used by the viewer.

There is an important caveat: this is partly circular because the combined
score already contains these metrics. Shape IoU has the largest weight
\((0.40)\), so high agreement between shape IoU and the combined score is
expected. Therefore, this check should be interpreted as a redundancy and
stability check, not as proof that shape IoU is physically correct.

\subsection{Best-Target Agreement}

For each source sheet \(s\), I compared:

\[
\arg\max_t C(s,t)
\]

against:

\[
\arg\max_t m(s,t),
\]

where \(m\) is one individual metric such as shape IoU or support Jaccard.

If the same target is selected, the individual metric agrees with the current
combined score for that source sheet.

The observed best-target agreement was:

\begin{center}
\begin{tabular}{lrrrrr}
\toprule
Dataset & Shape IoU & BBox IoU & Support Jaccard & Area Ratio & Centroid \\
\midrule
MVK\_s1 & 95.6\% & 94.2\% & 50.0\% & 65.9\% & 90.9\% \\
MVK\_s2 & 93.7\% & 91.4\% & 48.6\% & 58.1\% & 88.0\% \\
stilbene & 96.1\% & 89.3\% & 40.5\% & 47.7\% & 87.1\% \\
torusGrowing & 85.4\% & 44.2\% & 70.2\% & 43.1\% & 42.7\% \\
torusMoving & 37.7\% & 29.2\% & 46.5\% & 89.2\% & 31.2\% \\
\bottomrule
\end{tabular}
\end{center}

For the real datasets, shape IoU almost always selects the same target as the
combined score. This suggests that shape IoU is the dominant useful range-shape
signal. Bounding-box IoU is also useful, but coarser. Support Jaccard often
chooses different targets, which means it captures a different notion of
similarity.

\subsection{Domain-vs-Shape Disagreement}

I also compared the best target selected by vertex/domain overlap against the
best target selected by the shape combined score. The agreement was low:

\begin{center}
\begin{tabular}{lr}
\toprule
Dataset & Agreement \\
\midrule
MVK\_s1 & 9.7\% \\
MVK\_s2 & 9.6\% \\
stilbene & 8.7\% \\
torusGrowing & 16.0\% \\
torusMoving & 1.8\% \\
\bottomrule
\end{tabular}
\end{center}

This is not necessarily a failure. It means domain overlap and range-shape
overlap answer different questions:

\begin{itemize}[leftmargin=*]
\item Domain overlap asks whether the same spatial grid vertices participate.
\item Shape overlap asks whether the sheet occupies a similar region in the
  bivariate range space.
\end{itemize}

For a paper, this disagreement can be useful because it shows that the measures
are complementary rather than redundant.

\section{What Is the Continuity Threshold?}

The continuity threshold used in the previous summary is not currently part of
the main pipeline. It was an exploratory post-processing heuristic used to find
candidate event intervals.

For adjacent timesteps \(T_i\) and \(T_{i+1}\), let \(C(s,t)\) be the combined
shape score between source sheet \(s \in T_i\) and target sheet \(t \in
T_{i+1}\).

For a source sheet \(s\), define:

\[
\operatorname{bestOut}(s) = \max_t C(s,t).
\]

For a target sheet \(t\), define:

\[
\operatorname{bestIn}(t) = \max_s C(s,t).
\]

Using a threshold \(\theta = 0.5\):

\begin{itemize}[leftmargin=*]
\item if \(\operatorname{bestOut}(s) < \theta\), source sheet \(s\) is counted
  as having no strong continuation;
\item if \(\operatorname{bestIn}(t) < \theta\), target sheet \(t\) is counted
  as having no strong predecessor;
\item if one source sheet has two or more targets with \(C(s,t) \ge \theta\),
  it is counted as a possible split;
\item if one target sheet has two or more sources with \(C(s,t) \ge \theta\),
  it is counted as a possible merge.
\end{itemize}

The threshold \(\theta=0.5\) was chosen only as a mid-scale heuristic for a
score in \([0,1]\). It is not calibrated. Before using this in the paper, the
pipeline should compute sensitivity over several thresholds, for example:

\[
\theta \in \{0.3,0.4,0.5,0.6,0.7\}.
\]

If the same important intervals remain important across this range, the event
detection is more convincing. If the intervals change completely, the threshold
is too arbitrary.

\section{Current Candidate Intervals}

These are candidate intervals found from the current generated data. They
should be treated as hypotheses for closer inspection, not final paper claims.

\subsection{MVK\_s1}

The strongest candidate event region is:

\begin{verbatim}
step_01180 -> step_01220
\end{verbatim}

The most notable adjacent pair is:

\begin{verbatim}
step_01200 -> step_01220
\end{verbatim}

For this pair, the mean best combined score is approximately \(0.396\), and
18 out of the 20 top sheets have best continuation below the exploratory
threshold \(0.5\).

There is also one strong persistent feature: the rank-1 sheet persists across
all 82 available timesteps.

\subsection{MVK\_s2}

The strongest candidate event region is:

\begin{verbatim}
step_01120 -> step_01180
\end{verbatim}

Notable adjacent pairs are:

\begin{verbatim}
step_01140 -> step_01160
step_01160 -> step_01180
\end{verbatim}

The dominant feature appears to reconfigure in this interval. The top-sheet
area also changes sharply around this region.

\subsection{stilbene}

The strongest candidate event region is:

\begin{verbatim}
step_11910 -> step_11925
\end{verbatim}

The most notable adjacent pairs are:

\begin{verbatim}
step_11915 -> step_11920
step_11920 -> step_11925
\end{verbatim}

For both pairs, all 20 top sheets had best continuation below the exploratory
threshold \(0.5\). This makes the region a strong candidate for a paper figure.

Other large top-area jumps were found at:

\begin{verbatim}
step_01220 -> step_01240
step_06360 -> step_06380
step_06440 -> step_06480
step_07660 -> step_07680
step_09580 -> step_09600
\end{verbatim}

These need visual and fiber-surface inspection before interpretation.

\subsection{Synthetic Torus}

The synthetic torus datasets are useful for validation, but the current results
are not yet clean enough for a main paper claim.

For \texttt{torusGrowing}, rank-1, rank-2, and rank-3 tracks persist across all
25 timesteps. However, many tiny or zero-vertex sheets appear among the top 20,
which makes birth/death counts noisy.

For \texttt{torusMoving}, the rank-1 track persists across all 25 timesteps,
but several area jumps appear. These may reflect synthetic construction
artifacts rather than meaningful Reeb-space events.


\section{Implemented Analysis Stage}

A first reproducible analysis stage has now been implemented as:

\begin{verbatim}
stage_06_analyze_tracking_results.py
\end{verbatim}

It consumes the existing viewer result file:

\begin{verbatim}
sankey/unified_sankey_viewer/data.json
\end{verbatim}

It does not recompute Reeb spaces, sheet geometry, sheet images, or
fiber-surface images. By default it analyzes \texttt{common.BASE\_DIR}:

\begin{verbatim}
python3 stage_06_analyze_tracking_results.py
\end{verbatim}

It can also regenerate analysis outputs for all known datasets:

\begin{verbatim}
python3 stage_06_analyze_tracking_results.py --all-known
\end{verbatim}

For each dataset, outputs are written to:

\begin{verbatim}
<dataset>/sankey/tracking_analysis/
\end{verbatim}

The current outputs are:

\begin{verbatim}
metric_summary.csv
best_target_agreement.csv
event_scores.csv
sheet_lifetimes.csv
interesting_intervals.json
analysis_metadata.json
plots/*.png
plots/*.pdf
\end{verbatim}

The script also supports explicit threshold control, for example:

\begin{verbatim}
python3 stage_06_analyze_tracking_results.py --thresholds 0.3,0.4,0.5,0.6,0.7
\end{verbatim}

The pipeline has an optional flag, \texttt{RUN\_STAGE\_6\_TRACKING\_ANALYSIS},
which runs this stage after the unified Sankey viewer when enabled.


\section{Remaining Work After the Analysis Stage}

The first reproducible analysis stage is now implemented. The remaining work is
therefore not to manually recompute the diagnostics, but to validate and
interpret the generated outputs.

The next concrete steps are:

\begin{enumerate}[leftmargin=*]
\item Inspect \texttt{interesting\_intervals.json} for each dataset and select a
  small number of intervals for paper figures.
\item Use \texttt{event\_scores.csv} to check threshold sensitivity across
  \(0.3,0.4,0.5,0.6,0.7\). A convincing event should remain visible across
  nearby thresholds.
\item Use \texttt{sheet\_lifetimes.csv} to identify persistent tracks and decide
  which tracks should be shown or summarized.
\item Use \texttt{best\_target\_agreement.csv} and
  \texttt{metric\_summary.csv} to justify which metrics are primary and which
  should be treated as supplementary.
\item Manually inspect sheet and fiber-surface images linked from
  \texttt{interesting\_intervals.json} before making chemical or structural
  claims.
\item Decide how to report zero-regular-vertex sheets, especially for the
  synthetic torus datasets.
\end{enumerate}

The implemented event counts still depend on heuristic thresholds. Therefore,
threshold sensitivity and manual visual inspection should be treated as part of
the paper result-selection process.

\section{How This Fits the Paper}

The likely paper story is:

\begin{enumerate}[leftmargin=*]
\item Reeb spaces decompose each timestep of a bivariate field into sheets.
\item These sheets can be treated as time-varying bivariate features.
\item Tracking sheets across timesteps reveals persistent features,
  reconfigurations, and birth/death-like behavior.
\item Domain overlap and range-shape similarity are complementary. They should
  not be collapsed into a single story too early.
\item Shape IoU appears to be the most promising primary range-shape measure,
  while support Jaccard, area ratio, bounding-box IoU, and centroid similarity
  are useful as diagnostics or supplementary measures.
\end{enumerate}

Compared with the earlier continuous-scatterplot and image-moment paper, this
approach provides a feature-level view. The earlier method summarized atom- or
segment-level CSP behavior using moment tracks. The Reeb-space tracking
approach instead tracks sheet features directly and can show how those features
persist, split, merge, appear, or disappear.

\section{Cautions Before Making Paper Claims}

The current results are promising, but the following should be done before
writing final claims:

\begin{itemize}[leftmargin=*]
\item implement threshold sensitivity analysis;
\item manually inspect candidate intervals in the viewer;
\item check corresponding sheet and fiber-surface images;
\item decide how to handle zero-regular-vertex sheets in paper summaries;
\item clean or clearly label the synthetic torus results;
\item avoid claiming chemical interpretation until the candidate intervals are
  compared against domain knowledge and the prior CSP/moment findings.
\end{itemize}

\section{References for Context}

The previous CSP and image-moment paper used for context is:

\begin{center}
\url{https://vgl.csa.iisc.ac.in/pdf/pub/Paper_Continuous_Scatterplot_and_Image_Moments_for_Time_Varying_Bivariate_Field_Analysis.pdf}
\end{center}

The supplementary material is:

\begin{center}
\url{https://vgl.csa.iisc.ac.in/pdf/pub/Supp_Material_Continuous_Scatterplot_and_Image_Moments_for_Time_Varying_Bivariate_Field_Analysis.pdf}
\end{center}

\end{document}
